%!TEX program = lualatex
\documentclass[ngerman,headheight=70pt]{scrartcl}
\usepackage[ngerman]{babel}
\usepackage{blindtext}
\usepackage[utf8]{luainputenc}
\usepackage{scrlayer-scrpage}
\usepackage[T1]{fontenc}
\usepackage{fontspec}
\usepackage{multicol}
\usepackage{enumitem}
\usepackage[german=quotes]{csquotes}
\usepackage[bottom=2cm,footskip=8mm]{geometry}
\setmainfont{Source Sans Pro}
\MakeOuterQuote{"}

\setlength{\parindent}{0pt}
\setlength{\parskip}{6pt}

\pagestyle{plain.scrheadings}
\thispagestyle{scrheadings}
\lohead{\vspace*{-1.5cm}\includegraphics[scale=0.4]{../images/stupa_logo_rot.pdf}}
\rohead{\\Präsidium des\\Studierendenparlaments}

\lofoot[Universität Hamburg · Präsidium des Studierendenparlaments\\\hrule Von-Melle-Park 5 · D-20146 Hamburg]{Universität Hamburg · Präsidium des Studierendenparlaments\\\hrule Von-Melle-Park 5 · D-20146 Hamburg}
\cofoot[]{}
\cohead[\thepage]{}

\renewcommand{\thesubsection}{TOP \arabic{subsection}}
\renewcommand{\thesubsubsection}{\arabic{subsubsection}.}
\setcounter{subsection}{-1}

\begin{document}
    UHH · StuPa-Präsidium · Von-Melle-Park 5 · D-20146 Hamburg

    \section*{Beschluss des Studierendenparlamentes vom 14. April 2016}
    \subsection*{Gegen Rechts: Aufklärung, Kultur, Solidarität!}

	Vor über 160 Jahren bereits war erkennbar: Not, Angst und Gewalt sind
    menschengemacht und überwindbar. Die dramatische weltweite Ungleichverteilung
    von Vermögen und Einkommen spitzt sich mit rasanter Geschwindigkeit weiter zu.
    Produktivität und Reichtum aber könnten längst allen Menschen ein Leben in
    Frieden und Würde ermöglichen. Die soziale Angleichung der Lebensverhältnisse ist
	daher eine globale Notwendigkeit: für Frieden, Menschenrecht und Demokratie.

    Gegen diese – bereits der seit Befreiung von Faschismus und Krieg 1945 – mit
    der Gründung der Vereinten Nationen weltweit verbreitete Einsicht machen
    wieder rechte Hetzer mobil: Mit "Pegida" (usw.) und AfD formieren sie sich neu.
    Sie propagieren die Ungleichheit der Menschen mit Rassismus, Nationalismus und
    Sozialdarwinismus; sie wenden sich gegen Arbeitnehmerrechte, sozialstaatliche
    Errungenschaften, Grundrechte, gegen eine aufgeklärte Pluralität der Lebensweisen
    und forcieren die Konkurrenz bis zum Aufruf zu offener Gewalt. Sie und
    ihresgleichen in anderen Formationen der politischen Rechten sind die größten
    Scharfmacher für Krieg und Militarismus. Ihr Programm ist menschenfeindlich,
    unsozial und unrealistisch. Diese neu-alte Zusammenfassung von Autoritarismus,
    Spießertum und Profitkonformität ist eine irrationale Reaktion auf die massive
    Prekarisierung und Vereinzelung der Lebensverhältnisse in Folge neoliberaler
    "Reformen" und auf die breit getragene Flüchtlingssolidarität und die in der
    Bevölkerung wachsende Kritik an unsozialer und kriegerischer Politik des
    Establishments.

    Die Studierendenschaft positioniert sich als Akteurin des Umbruchs aus einer
    Phase extremer globaler Verschärfung von Ausbeutung und Machtpolitik hin zum
    solidarischen Eintreten für globale Gerechtigkeit. Sie ist in diesem
    Bewusstsein auch solidarisch mit allen Geflüchteten. Die Studierendenschaft
    will aktiver Teil einer weltoffenen solidarischen Kulturbildung sein. Sie sieht
    sich deshalb in der Pflicht, ihre gesellschaftliche Aufklärungsarbeit gegen
    Rechts, ihr Engagement für Frieden, Abrüstung, sozialen Fortschritt und für
    umfassende Demokratie zusammenzuführen und zu intensivieren. Das
    Studierendenparlament will diese Orientierung forcieren.

    \begin{enumerate}
        \item Das Studierendenparlament beauftragt seinen Ausschuss gegen Rechts
        damit, im Sommersemester einen Aktionstag gegen Rechts durchzuführen.
        Dieser soll sich besonders der analytischen und kulturkritischen
        Auseinandersetzung mit der AfD widmen und über eine wirkliche, soziale und
        (humorvoll) aufgeklärte Alternative zum neoliberalen Gesellschaftsmodell
        informieren.
        \item Das Studierendenparlament beauftragt den AStA in Verständigung mit
        dem Ausschuss gegen Rechts, die Mobilisierung zu der Veranstaltung
        "Bücherverbrennung – Nie wieder!" (10. Mai 2016) zu verbreitern. Dabei
        soll auch ein Aufruf zu der Veranstaltung zur Verabschiedung in der
        nächsten StuPa-Sitzung vorbereitet werden.
    \end{enumerate}
\end{document}
